% Generato da build_cv.py — non modificare a mano
\begin{rubric}{Other activities}
\subrubric{Organization}
\entry*[Jul 2025 -- Jan 2026]
  \textbf{Member of the local Organizing Committee of the 4th Workshop of UMI Group ``Mathematics for Artificial Intelligence and Machine Learning''}\par Department of Mathematics G. Castelnuovo, Sapienza Università di Roma\par Fourth edition of the workshop of the UMI Group “Mathematics for Artificial Intelligence and Machine Learning”, following the previous ones held in Milan, Turin, and Bari.
\entry*[Nov 2024 -- Sep 2025]
  \textbf{Co-organizer of the Summer School ``Mathematical methods for high-dimensional data''}\par Department of Mathematics G. Castelnuovo, Sapienza Università di Roma\par The school will open the thematic period on Data Science of the ``Excellence program'' for the Department of Mathematics: see the \href{https://sites.google.com/view/math-high-dimensional-data/home}{website.}
\entry*[2024 -- today]
  \textbf{Responsible for communication}\par Communication roles for the FAIR foundation, Spoke 5 (WP5.5) for the organization of scientific and dissemination events.
\entry*[2024]
  \textbf{Organizer of scientific meetings}\par I covered organizing roles for the FAIR Spoke 5 scientific meeting, consisting in discussions and crossfertilization between working packages for the FAIR (Future Artificial Intelligence Research).
\subrubric{Other roles}
\entry*[2022]
  \textbf{Project reviewer}\par I have also been a project reviewer for the ISF
\entry*[2017 -- today]
  \textbf{Referee for high-impact journals}\par Among others: \emph{Physical Reviews X; IEEE Transaction on Neural Networks and Learning Systems; Nature Scientific Reports; Journal of Mathematical Physics; Journal of Physics A: Mathematical and Theoretical; Neural Computation; Journal of Physics: Complexity; Physica A: Statistical Mechanics and its Applications; Journal of Statistical Mechanics: Theory and Experiment; Machine Learning: Science and Technology; Physical Review E}
\subrubric{Visiting periods}
\entry*[16 -- 23 Jan 2023]
  \textbf{Alan Turing Institute, London (UK)}\par During the period, I attended the workshop ``Physics-informed Machine-Learning'', in which I took part to stimulating discussions about statistical mechanics of spin glasses and their application to Artificial Intelligence. I also gave a talk at the workshop about Dreaming Neural Networks and the possibility to extended the theory to Boltzmann Machines.
\entry*[1 -- 5 Oct 2018]
  \textbf{Department of Mathematics, King's College London, London (UK)}\par Ref: Prof. Alessia Annibale\par During the period, we discussed several topics in statistical mechanics of spin glasses, especially concerning the application of Machine Retrieval models to the theoretical modeling of biological processes. In this period, I also gave an invited talk. Finally, I also had the chance to discuss with brilliant researchers and professors, in particular with Prof. Reimer Kuhn.
\entry*[2018 (often)]
  \textbf{Department of Mathematics ``Guido Castelnuovo'', Sapienza Università di Roma, Rome (IT)}\par Ref: Prof. Elena Agliari\par During the period, we investigated unlearning theory for Machine Retrieval models. The discussions were crucial for the ideation and the theoretical investigation about Dreaming Neural Networks.
\entry*[2016 -- 2017 (often)]
  \textbf{Department of Physics ``Augusto Righi'', Alma Mater Studiorum, Bologna (IT)}\par Ref: Prof. Davide Fioravanti\par During the period, we investigated the relation between supersymmetric gauge theory on deformed background with quantization of classical integrable systems.
\end{rubric}
